\subsection{正则扩张的刻画}

命题 9. --- 设 K 是一个域， \(\bar{K}\) 是 K 的一个代数闭包，L 是 K 的一个扩张。则下列条件等价：

\begin{enumerate}
\def\labelenumi{\alph{enumi})}
\item
  L 在 K 上是可分的，且 K 在 L 中相对代数封闭。
\item
  L 是 K 的正则扩张。
\item
  环 \(\bar{\mathbf{K}}\otimes_{\mathbb{K}}\mathbf{L}\) 是整环。
\item
  设 \(\bar{L}\) 是 L 的一个代数闭包。则 L 与 K 在 \(\bar{L}\) 中的相对代数闭包在 K 上线性不相交。
\end{enumerate}

此外，当这些条件成立时，则 \(\bar{\mathbf{K}}\otimes_{\mathbf{K}}\mathbf{L}\) 是一个域。

\(a) \Rightarrow b)\) ：设 M 是 K 的一个扩张。在 \(a\) ) 的假设下，环 M \(\otimes_{K} L\) 是整环（V，第 140 页，推论）。

\(b) \Rightarrow c)\) ：这由定义 2 得出。

\(c) \Rightarrow d): \text{记号为 } d) \text{ 我们可以识别 } \bar{\mathbf{K}} \text{ 取其相对代数闭包 } \mathbf{K} \text{ 在 } \bar{\mathbf{L}} (\mathbf{V}, \text{ p. 22, Ex. 2}). \text{ 设环 } \mathbf{A} = \bar{\mathbf{K}} \otimes_{\mathbf{K}} \mathbf{L} \text{ 是整环。令 } \mathbf{E} \text{ 为其子扩张 } \bar{\mathbf{K}} \text{ 有限次数扩张 } \mathbf{K}; \text{ 子环}\)

\(E \otimes_{K} L\) 这个 A 的子环是整环，从而是 L 上的一个有限次代数；由 V，第 10 页的推论，它是一个域。由于 \(\bar{K}\) 是上述类型的扩张 E 组成的递增有向集的并，A 是一个域（V，第 11 页，命题 3）。于是，A 到 \(\bar{L}\) 中的、把 \(x \otimes y\) 映为 xy（对于 \(x \in \bar{K}\) 和 \(y \in L\) ）的典范同态是单射，所以 L 和 \(\bar{K}\) 在 K 上线性不相交。

\(d) \Rightarrow a): \text{在其假设下 } d) \text{ 有 } L \cap \bar{K} = K, \text{ 所以 } K \text{ 在其中相对代数闭 } L; \text{ 此外，若 } p \text{ 是其特征指数 } K, \text{ 域 } L \text{ 与其线性不交 } K^{p^{-\infty}} \text{ 在其上 } K, \text{ 因此 } L \text{ 在其上可分离 } K (V, p. 123, Cor. 1).\)

推论 1. --- 设 A 是域 K 上的一个代数。要使 A 是正则 K-代数，必要且充分条件是环 \(\bar{\mathbf{K}}\otimes_{\mathbf{K}}\mathbf{A}\) 是整环。

所述条件显然是必要的。反之，假设 \(\bar{K} \otimes_{K} A\) 是一个整环，并记 E 为 A 的分式域。由命题 4（V，第 141 页），环 \(\bar{K} \otimes_{K} E\) 是一个整环，因此由命题 9，E 是 K 的正则扩张；由 V，第 141 页的推论，我们得出 A 作为 K-代数是正则的。

推论 2. --- 设 K 是一个代数闭域。每个作为整环的 K-代数都是正则 K-代数。特别地，K 的每个扩张都是正则的。

这由推论 1 得出。

推论 3. --- 设 K 是一个代数闭域。若 A 和 B 是两个作为整环的 K-代数，则 A ⊗\_K B 亦然。

根据推论 2，A 和 B 是正则 K-代数，只需应用命题 2（V，第 141 页）。

